% !TEX program = pdflatex
\documentclass[handout,aspectratio=169,11pt]{beamer}

% =========================
% PAQUETES
% =========================
\usepackage[spanish]{babel}
\usepackage[utf8]{inputenc}
\usepackage[T1]{fontenc}
\usepackage{lmodern}
\usepackage{amsmath,amsfonts,amssymb}
\usepackage{xcolor}

% =========================
% COLORES UST (PRO)
% =========================
\definecolor{USTBlue}{HTML}{003A8F}
\definecolor{USTLight}{HTML}{E6ECF5}
\definecolor{USTAccent}{HTML}{4CAF50} % verde agronomía

\setbeamercolor{title}{fg=white,bg=USTBlue}
\setbeamercolor{frametitle}{fg=white,bg=USTBlue}
\setbeamercolor{structure}{fg=USTBlue}
\setbeamercolor{block title}{bg=USTBlue, fg=white}
\setbeamercolor{block body}{bg=USTLight}

% =========================
% PIE DE PÁGINA
% =========================
\setbeamertemplate{footline}{
	\leavevmode%
	\hbox{
		\begin{beamercolorbox}[wd=.7\paperwidth,ht=2.5ex,dp=1ex,left]{author in head/foot}
			\hspace{0.3cm}MAT-016 | Expresiones Algebraicas
		\end{beamercolorbox}
		\begin{beamercolorbox}[wd=.3\paperwidth,ht=2.5ex,dp=1ex,right]{date in head/foot}
			\insertframenumber{} / \inserttotalframenumber \hspace{0.3cm}
	\end{beamercolorbox}}
}

\setbeamertemplate{navigation symbols}{}

% =========================
% PORTADA
% =========================
\title{Guía de Ejercicios 09}
\subtitle{Expresiones Algebraicas en Agronomía}
\author{Profesor Luis Tulio González Alcaino}
\institute{Universidad Santo Tomás \\ Departamento de Ciencias Básicas}
\date{Semestre I - 2026}

% =========================
\begin{document}
	
	% -----------------------------------
	\begin{frame}
		\titlepage
	\end{frame}
	
	% -----------------------------------
	\begin{frame}{Objetivo de la guía}
		\begin{itemize}
			\item Reducir expresiones algebraicas
			\item Identificar términos semejantes
			\item Aplicar operaciones con polinomios
			\item Resolver problemas en contexto agrícola
		\end{itemize}
	\end{frame}
	
	% ===================================
	\section{Reducción de expresiones}
	
	% -----------------------------------
	\begin{frame}{Actividad 1}
		\small
		Simplifique y determine el grado:
		
		\begin{enumerate}
			\item $5ab^2 - 3ab + 15b^2a - 12ab + 7a^2b$
			\item $18m^5n^3r - 7n^5m^3r + (3m^5n^3r - 12)$
			\item $14 + [12xy - (3x^2y^2 - 4xy) + 8xy + 10]$
			\item $-10t + [-( -2t + 3) + 12] - (14 + 2t + 11t)$
		\end{enumerate}
	\end{frame}
	
	% -----------------------------------
	\begin{frame}{Actividad 2}
		\small
		Contexto agrícola:
		
		\begin{enumerate}
			\item $x^2yz - (18xy^2z + 10xyz^2)$
			\item $11a^2b - 32a^2b + 14ab^2 - 3a^2b$
			\item $10x^2yz - 12x^2yz - 8yx^2z + 42yx^2z$
			\item $120mn - 54mn + 100m - 30mn - 28m$
		\end{enumerate}
	\end{frame}
	
	% -----------------------------------
	\begin{frame}{Actividad 3}
		\small
		\begin{enumerate}
			\item $35pq + 30 + 78pq - 60pq - 13$
			\item $54rt^2 - (12rt^2 - 5rt) + (21rt^2 - 2rt)$
		\end{enumerate}
	\end{frame}
	
	% ===================================
	\section{Operaciones con polinomios}
	
	% -----------------------------------
	\begin{frame}{Modelos costo agrícola}
		\small
		
		\textbf{Costo Producción:} $A = 3x^2 - 5xy - 2xy^2$
		
		\textbf{Costo Fertilización:} $B = 10xy + 7x^2 - 3y^2$
		
		\textbf{Costo Plagas:} $C = -4y^2 - 8xy + 10$
		
	\end{frame}
	
	% -----------------------------------
	\begin{frame}{Ejercicios}
		\small
		\begin{enumerate}
			\item $A - C + 34$
			\item $-A + B - D$
			\item $C - (B + A)$
			\item $-(A + D) - 7$
			\item $2C - [4B - (3B - C) - 2D] - D$
		\end{enumerate}
	\end{frame}
	
	% -----------------------------------
	\begin{frame}{Crecimiento vegetal}
		\small
		
		$A = 2x^3 - x^2 + 4x - 5$
		
		$B = x^3 - 8x - 2$
		
		$C = x^2 - 4x - 2$
		
	\end{frame}
	
	% -----------------------------------
	\begin{frame}{Ejercicios}
		\small
		\begin{enumerate}
			\item $A - (B - C)$
			\item $A - 2B$
			\item $A - B + 3C$
			\item $3(B - 2A)$
			\item $4A - [2B + C - (3B - 2A)] - A$
		\end{enumerate}
	\end{frame}
	
	% ===================================
	\section{Problemas aplicados}
	
	% -----------------------------------
	\begin{frame}{Problema 1: Fertilización}
		En un cultivo:
		
		\begin{itemize}
			\item Base: $(2n + 15)$
			\item Refuerzo: $(3n - 10)$
		\end{itemize}
		
		\pause
		
		\textbf{Pregunta:}  
		¿Cuál es la cantidad total aplicada?
		
	\end{frame}
	
	% -----------------------------------
	\begin{frame}{Problema 2: Nutrientes}
		\small
		\begin{itemize}
			\item Suelo: $45x^2 - 12xy + 30y^2$
			\item Fertilizante: $-7x^2 + 24xy - 3y^2$
			\item Abono: $30x^2 - 10xy + y^2$
		\end{itemize}
		
		\pause
		
		\textbf{Preguntas:}
		\begin{itemize}
			\item Sin abono
			\item Con abono
		\end{itemize}
		
	\end{frame}
	
	% -----------------------------------
	\begin{frame}{Problema 3: Pesticidas}
		\small
		Disponible:
		\[
		50a + 20b + 30
		\]
		
		Aplicación:
		
		\begin{itemize}
			\item Parcela 1: $5a - 35b + 42$
			\item Parcela 2: $32a + 10b - 22$
		\end{itemize}
		
		\pause
		
		\textbf{Pregunta:}  
		¿Cuánto queda?
		
	\end{frame}
	
	% -----------------------------------
	\begin{frame}{Cierre}
		\begin{block}{Idea clave}
			Reducir expresiones algebraicas permite modelar situaciones reales en la agricultura.
		\end{block}
		
		\pause
		
		\begin{block}{Aplicación}
			Fertilización, riego, crecimiento y costos pueden representarse algebraicamente.
		\end{block}
		
	\end{frame}
	
	% =========================
\end{document}